\chapter{Saúde e autocuidado mental no mundo digital}
% TODO: Move these refs into bibtex
No quesito de saúde mental, a bipolaridade ou pelo seu nome científico,
Transtorno Afetivo Bipolar (TAB) é uma das doenças mentais mais comuns e
responsável por uma gama muito grande de disfunções mentais. De acordo com os
dados da Organização Mundial da Saúde (WHO, 2017), afeta 0,5\% da população
mundial. No contexto do Brasil, a bipolaridade, segundo os dados da Associação
Brasileira de Psiquiatria, estima que a bipolaridade aflige cerca de 2\% da
população, equivalente à uns 4,26 milhões de pessoas. Esse número, embora soe
pouco expressivo, levando em consideração o contingente de pessoas no país, é
preocupante, dado o fato de que a bipolaridade é (1) complexa de se detectar e
(2) muitas vezes, acaba passando despercebida ou enquadrada dentro de um quadro
de depressão unipolar (depressão).

Ademais, portadores de TAB estão fortemente associados a comorbidades e
hospitalizações, com cerca de 15\% dos portadores de TAB cometendo suicídio
(JUDD et al. 2002). Com esse recorte em mãos, os principais tratamentos, para
além da medicação adequada, são atividades físicas/exercícios regulares,
rotinas estruturadas de sono, acompanhamento psicoterapêutico, e, sobretudo,
monitoramento constante do humor. O auto registro diário é reconhecido como
ferramenta essencial para antecipar episódios maníaco/depressivo, facilitando
intervenções imediatas antes da crise plena~\cite{Faurholt:2016}. 

Nesse contexto, o uso de tecnologias digitais tem se mostrado uma via
promissora para sistematizar esse tipo de acompanhamento. Pesquisas indicam que
pessoas com transtornos do espectro bipolar demonstram maior adesão ao
monitoramento emocional quando este é feito de forma acessível, visual e
integrada ao seu cotidiano~\cite{Iwaya:22}.

Historicamente, a prática da escrita reflexiva tem sido aplicada como mecanismo
de autocuidado. As suas raízes são antigas e permanentes, de forma que
filósofos gregos, como Sêneca e Marcus Aurelius mantinham diários com o
objetivo de revisar ações e promover uma espécie de meditação moral e
autotransformação, para além de servir como um espaço de notas pessoais. Já, na
era moderna, com o desenvolvimento da psicanálise e uma melhoria na compreensão
afetiva, vemos surgir técnicas como a da journal therapy (Progoff, 1996),
mecanismo que sistematiza o uso da escrita como forma de promover bem estar
emocional. Formas de auto registro ganham ainda mais um fator elementar dentro
com o desenvolvimento da psicologia cognitivo-comportamental (TCC), da qual
aprofunda esse elemento como ferramenta para identificação de padrões de
pensamentos disfuncionais (BECK, 1997).

Trabalhando ainda com esse fator, o surgimento e desenvolvimento da internet, e
portanto, o acesso a mesma, inicia-se uma progressiva digitalização do que
antes era analógico. Processos que eram previamente feitos na caneta e no papel
passam a ser manualmente migrados para suas respectivas versões eletrônicas.
Atividades como escrever cartas, manter álbuns de fotografias, consultar
enciclopédias passam por uma rápida transformação ao longo das últimas três
décadas. O que antes precisava de um suporte físico é transportado para dentro
de computadores, tablets e mais recentemente, para os smartphones. Essa
transformação se insere no que Castells (1999) chama de “sociedade em rede”,
onde o digital se torna a principal estrutura por onde circula o conhecimento,
a comunicação e a cultura, fomentando uma interconexão digital de saberes e
práticas cotidianas.

Um exemplo marcante que pode ser citado, é a substituição de bibliotecas
físicas por acervos digitais, como por exemplo a Biblioteca Nacional Digital do
Brasil, que oferece acesso remoto a obras históricas e documentos raros
(BIBLIOTECA NACIONAL, 2024). Outro exemplo notável, foi a mudança nos sistemas
de ensino, que incorporaram ambientes virtuais de aprendizagem, inclusive
significativamente acelerados pela pandemia da COVID-19, transformando práticas
educacionais em escala global (HODGES et al., 2020). Na esfera da comunicação
interpessoal, as cartas deram lugar ao e-mail, e posteriormente, a aplicativos
de comunicação instantânea como o WhatsApp (LÉVY, 1999; CASTELLS, 2003).

No âmago do espaço pessoal, planejamento do cotidiano, que envolvia agendas de
papel, listas de tarefas, cadernos de anotações, foram absorvidas por
ferramentas digitais, como calendários online e aplicativos de produtividade
(TURKLE, 2011). Esse processo demonstra que a digitalização, não se mantém
apenas como um processo técnico e mecânico, mas como um elemento cultural de
adaptação e reconfiguração das práticas humanas, que caminha para a
virtualização das experiências (MILLER; HORST, 2012). É nesse contexto, onde o
cotidiano é constantemente traduzido em dados, que mesmo as informações
subjetivas, como sentimentos e estados emocionais, que se centraliza a
transição dos diários pessoais para os registros digitais simples, e
posteriormente, de forma estruturada, com o surgimento os e-journals ou diários
digitais marca a virtualização da subjetividade (MOROZOV, 2013; LUPTON, 2016).
